\documentclass[12pt]{article}
\usepackage[T1]{fontenc}
\usepackage[utf8]{inputenc}
\usepackage{lmodern}
\usepackage{textcomp}
\usepackage{enumitem}
\usepackage{geometry}
\usepackage{graphicx}
\usepackage{amsmath, amssymb}
\geometry{margin=1in}
\begin{document}
\begin{center}
Political Science - Undergraduate Level Assessment 13 \
Total Marks: 40 \
Answer all questions.
\end{center}

\section*{Multiple Choice (10 marks)}
\begin{enumerate}[label=\arabic*.]
\item Which of the following is true of the Supreme Court?
\begin{enumerate}[label=(\alph*)]
    \item Every case appealed to the Supreme Court is ruled upon by the court.
    \item The court helps set the public agenda by deciding which appeals to hear.
    \item The court hears all cases when two or more justices agree that the case has merit.
    \item In deciding cases, the chief justice's vote counts as two votes.
\end{enumerate}
\item What was the key difference between US expansion pre- and post- 1865?
\begin{enumerate}[label=(\alph*)]
    \item US expansion was based on territory rather than markets post-1865
    \item US expansion was based on markets rather than territory post-1865
    \item US expansion was limited to Latin America post-1865
    \item US expansion ended after 1865
\end{enumerate}
\item What, according to Systemic theories, is the primary determinant of a state's foreign policy?
\begin{enumerate}[label=(\alph*)]
    \item The character of a state's leader
    \item The distribution of power in the international system
    \item The distribution of power within a state's governmental system
    \item A state's political ideology
\end{enumerate}
\item What was meant by the term 'New World Order'?
\begin{enumerate}[label=(\alph*)]
    \item A new democratic internationalism led by the United States
    \item A new balance of power between the US and China
    \item A new global economic framework
    \item A new era of globalization
\end{enumerate}
\item A constitutional amendment would be required to ban flag burning because that activity is currently protected by the right to
\begin{enumerate}[label=(\alph*)]
    \item due process
    \item assembly
    \item free exercise of religion
    \item free speech
\end{enumerate}
\end{enumerate}

\section*{True / False (10 marks)}
\textit{Answer True or False}
\begin{enumerate}[label=\arabic*.]
\setcounter{enumi}{5}
\item True or False: The oversight of executive branch agencies is an example of an implied power of Congress derived from its legislative authority.
\item True or False: Congress has the authority to override a presidential veto with a two-thirds majority vote in both chambers.
\item True or False: The Supreme Court has applied the abolition of slavery clause of the Fourteenth Amendment to enforce privacy rights against state law.
\item True or False: The Supreme Court grants writs of certiorari exclusively in cases involving constitutional issues.
\item True or False: The National Security Council is the main agency that conducts diplomatic negotiations for U.S. foreign policy.
\end{enumerate}

\section*{Long Form Response (20 marks)}
\textit{Answer each question with a clear and complete explanation.}
\begin{enumerate}[label=\arabic*.]
\setcounter{enumi}{10}
\item Discuss how the examination of state practice and policy since the Cold War supports the notion of a customary right to humanitarian intervention, while also considering the absence of an obligation.
\item Discuss how game theory and deterrence theory contributed to intellectual advancements in security studies during the 'golden age,' highlighting their impact on strategic concepts.
\end{enumerate}

\vfill
\noindent\textit{End of Paper}
\end{document}